\documentclass[12pt]{article}
\usepackage[utf8]{inputenc}
\usepackage[T1]{fontenc}
\usepackage{hyperref}
\title{Tennet Language Resources for K-12 Curriculum Development}
\author{Compiled by Claude Sonnet 5 (claude-sonnet-5) for Andr\'as Kornai}
\date{}
\begin{document}
\maketitle

\section{Introduction}
Tennet (ISO 639-3 tex; Glottocode tenn1246) is a Surmic language of Eastern Equatoria, South Sudan. Estimates of the number of speakers vary widely (4,000 to 30,000 depending on the source). General information: \url{https://en.wikipedia.org/wiki/Tennet_language}.

\section{Primary Educational Resources}
No primer, textbook or school materials were found in the searches. The Endangered Languages Project lists no standard orthography: \url{https://www.endangeredlanguages.com/lang/5060?hl=en}.

\section{Linguistic and Reference Materials}
Randall, Scott (1998), ``A Grammatical Sketch of Tennet'', in Dimmendaal and K\"oppe (eds.), \emph{Surmic Languages and Cultures} (bibliographic listing only): \url{https://endangeredlanguages.com/node/96782}.
Adelino Amargira's abstract on tone, which describes a tone-marking convention (acute for high, unmarked for low, circumflex for falling): \url{https://www.uni-koeln.de/phil-fak/afrikanistik/wocal/schedule/abstracts/8-8-3\%20Adelino\%20Amargira.pdf}.
Overview of the Surmic family: \url{https://en.wikipedia.org/wiki/Surmic_languages}.

\section{Religious and Community Resources}
None found.

\section{Digital and Audio Resources}
None found.

\section{Curriculum Development Context}
Documentation consists of a grammar sketch and a few papers on tone and morphology. There is no dictionary, reader or scripture that could be verified. Any curriculum would first need an agreed orthography, building on the tone-marking work above, and consultation with South Sudanese educators.
\end{document}